\section{Graham Schmidt}
% ============================================================
% Notes: Gram--Schmidt Orthogonalization (functions)
% Topic: Construct orthogonal / orthonormal set from a
%        linearly independent non-orthogonal set
% Source: Arfken & Weber (6th ed.), Ch. 10.3
% ============================================================

\section{Gram--Schmidt orthogonalization of functions}

\subsection{Goal and setting}
Given a set of \textbf{linearly independent} functions
\[
\{f_1(x), f_2(x), \dots, f_N(x)\}
\]
that are not necessarily orthogonal, we want to construct an \textbf{orthogonal}
(or orthonormal) set
\[
\{\phi_1(x), \phi_2(x), \dots, \phi_N(x)\}
\]
that spans the same subspace.

\subsection{Inner product and norm (weighted)}
In Sturm--Liouville theory the natural inner product is
\begin{equation}
\langle g|f\rangle := \int_a^b g^*(x)\,w(x)\,f(x)\,dx,
\qquad w(x)>0.
\end{equation}
The corresponding norm is
\begin{equation}
\|f\| := \sqrt{\langle f|f\rangle}.
\end{equation}

Orthogonality means
\begin{equation}
\langle \phi_m|\phi_n\rangle = 0,\qquad m\neq n.
\end{equation}
Orthonormality means
\begin{equation}
\langle \phi_m|\phi_n\rangle = \delta_{mn}.
\end{equation}

% ------------------------------------------------------------

\subsection{Gram--Schmidt construction (orthogonal set)}
Define recursively:
\begin{align}
u_1 &:= f_1,\\[4pt]
u_2 &:= f_2 - \frac{\langle u_1|f_2\rangle}{\langle u_1|u_1\rangle}\,u_1,\\[6pt]
u_3 &:= f_3
- \frac{\langle u_1|f_3\rangle}{\langle u_1|u_1\rangle}\,u_1
- \frac{\langle u_2|f_3\rangle}{\langle u_2|u_2\rangle}\,u_2,\\[6pt]
&\ \vdots\\[4pt]
u_n &:= f_n - \sum_{k=1}^{n-1}
\frac{\langle u_k|f_n\rangle}{\langle u_k|u_k\rangle}\,u_k.
\end{align}

Then the set $\{u_1,\dots,u_N\}$ is orthogonal:
\[
\langle u_m|u_n\rangle = 0 \quad (m\neq n).
\]

\subsection{Compact projection form}
Define the projection of $f$ onto $u_k$ as
\begin{equation}
\mathrm{proj}_{u_k}(f) :=
\frac{\langle u_k|f\rangle}{\langle u_k|u_k\rangle}\,u_k.
\end{equation}
Then Gram--Schmidt is simply
\begin{equation}
u_n = f_n - \sum_{k=1}^{n-1}\mathrm{proj}_{u_k}(f_n).
\end{equation}

% ------------------------------------------------------------

\subsection{Orthonormalization step}
To normalize each $u_n$, define
\begin{equation}
\phi_n := \frac{u_n}{\|u_n\|} =
\frac{u_n}{\sqrt{\langle u_n|u_n\rangle}}.
\end{equation}
Then
\[
\langle \phi_m|\phi_n\rangle = \delta_{mn}.
\]

% ------------------------------------------------------------

\subsection{Why Gram--Schmidt works (key identity)}
At step $n$, we subtract off all components of $f_n$ parallel to the already-built
orthogonal vectors $u_1,\dots,u_{n-1}$.
Thus $u_n$ lies in the span of $\{f_1,\dots,f_n\}$ but is orthogonal to the previous ones.

Specifically, for any $j<n$:
\begin{align}
\langle u_j|u_n\rangle
&= \left\langle u_j\left|f_n - \sum_{k=1}^{n-1}
\frac{\langle u_k|f_n\rangle}{\langle u_k|u_k\rangle}u_k\right.\right\rangle\\
&= \langle u_j|f_n\rangle - \sum_{k=1}^{n-1}
\frac{\langle u_k|f_n\rangle}{\langle u_k|u_k\rangle}\langle u_j|u_k\rangle\\
&= \langle u_j|f_n\rangle
- \frac{\langle u_j|f_n\rangle}{\langle u_j|u_j\rangle}\langle u_j|u_j\rangle\\
&=0.
\end{align}

% ------------------------------------------------------------

\subsection{Important practical notes}
\begin{itemize}
\item \textbf{Linear independence required:}
If $\{f_n\}$ are linearly dependent, then some $u_n$ becomes the zero function
($\|u_n\|=0$) and the algorithm breaks.
\item \textbf{Order matters:}
Different ordering of $\{f_n\}$ produces different orthogonal bases.
\item \textbf{Same subspace:}
For each $n$,
\[
\mathrm{span}\{u_1,\dots,u_n\} = \mathrm{span}\{f_1,\dots,f_n\}.
\]
\item \textbf{Degenerate eigenspaces:}
If an eigenvalue is degenerate, Gram--Schmidt produces an orthonormal basis
inside the degenerate subspace.
\end{itemize}

% ------------------------------------------------------------

\subsection{Two-function special case (fast formula)}
Given two linearly independent functions $f_1,f_2$:
\begin{align}
u_1 &= f_1,\\
u_2 &= f_2 - \frac{\langle f_1|f_2\rangle}{\langle f_1|f_1\rangle}\,f_1,\\
\phi_1 &= \frac{f_1}{\|f_1\|},\qquad
\phi_2 = \frac{u_2}{\|u_2\|}.
\end{align}

% ------------------------------------------------------------

\subsection{Example (typical weight)}
On $[-1,1]$ with weight $w(x)=1$:
\[
\langle g|f\rangle = \int_{-1}^{1} g(x)f(x)\,dx.
\]
Start with non-orthogonal polynomials:
\[
f_1(x)=1,\qquad f_2(x)=x,\qquad f_3(x)=x^2.
\]
Then
\[
u_1=1,
\qquad
u_2=x-\frac{\langle 1|x\rangle}{\langle 1|1\rangle}\cdot 1
= x,
\]
since $\int_{-1}^1 x\,dx=0$.
Next,
\[
u_3 = x^2 - \frac{\langle 1|x^2\rangle}{\langle 1|1\rangle}\cdot 1
      - \frac{\langle x|x^2\rangle}{\langle x|x\rangle}\cdot x.
\]
But $\int_{-1}^1 x^3\,dx=0$, so the $x$-projection vanishes and
\[
u_3 = x^2 - \frac{\int_{-1}^1 x^2 dx}{\int_{-1}^1 1\,dx}
= x^2 - \frac{2/3}{2}
= x^2 - \frac{1}{3}.
\]
These are proportional to the first Legendre polynomials:
\[
P_0=1,\quad P_1=x,\quad P_2=\frac{1}{2}(3x^2-1).
\]

% ------------------------------------------------------------

\subsection{Final high-yield summary}
\begin{itemize}
\item Gram--Schmidt turns $\{f_n\}$ into orthogonal $\{u_n\}$ by subtracting projections.
\item Weighted inner product:
\[
\langle g|f\rangle = \int_a^b g^*wf\,dx.
\]
\item Core recursion:
\[
u_n = f_n - \sum_{k=1}^{n-1}\frac{\langle u_k|f_n\rangle}{\langle u_k|u_k\rangle}u_k.
\]
\item Normalize:
\[
\phi_n = \frac{u_n}{\sqrt{\langle u_n|u_n\rangle}}.
\]
\item Used to build orthonormal bases in degenerate eigenspaces.
\end{itemize}
